% Pursuing Stacks (cotes 134-2 to 134-8): archive cote 134-2, site page 30 (PDF page 31).
% letter to Quillen. Draft compared: ps1.tex:771-835.
% AI-transcribed by Claude Opus 5.5 (claude-opus-5-5, adaptive thinking, effort high), Message Batches API, batch ps1 (msgbatch_01Rv5hCaNTtUt4JHx1yqoHPL), 2026-11-12; not yet checked by eye.
\scanpage{134-2/30}
\mspage{9}
\margin{\S 16}
that \uline{this ambiguity however is an irrelevant one.} To make this point
a little clearer, I could say that a third mathematician, informed of
the work of both, will readily think out a functor or rather a pair of
functors, associati\xout{m}ng to any structure of Mr. X one of Mr. Y and conver%
\hand{-}sely, in such a way that by composition of the two, we will associate
to a X-structure \hand{(}T say\hand{)} another T', which \deleted{of course} will not be isomor-%
phic to T of course, but endowed with a canonical \xout{xxxx} $\infty$-equivalence
(in the sense of Mr. X) $T \underset{\hand{\infty}}{\simeq} T'$ , and the same on the Mr. Y side. Most
probably, a fourth mathematician, faced with the same situation as the
third, will get \inserted{his own} \xout{\illegible{1}} pair of functors to reconciliate Mr X and Mr
Y , \inserted{which} \xout{\illegible{1}} very probably \xout{\unsure{his pair of functors}} won\hand{\unsure{'}}t be equivalent \xout{\illegible{1}} \hand{(}I mean
isomorphic\hand{)} to the previous one. Here however a fifth mathematician,
informed about this new perplexity, will probably show that the two
Y-structures U and U', associated by his two colleagues to \xout{the} an
X-structure T, while not isomorphic alas, admit however a canonical
$\infty$-equivalence between U and U'\inserted[typed]{(in the sense of the Y-theory)}. I could go on with a sixth mathematicia%
an, confronted with the same perplexity as the previous one, who winds
up with another \inserted{$\infty$-}equivalence between U and U' (without being informed of
the work of the fifth), and a seventh reconciling them by discovering
an $\infty$-equivalence between these equivalences. The story of course is
infinite, I better stop with seven mathematicians, a fair number now-%
adays to allow themselves getting involved with foundational matters ...
There should be a mathematical statement though resuming in finite terms
this infinite story, but in order to write it down I guess a minimum
amount of conceptual work, in the context of a given notion of $\infty$-grou-%
poids satisfying the desiderata (A) and (B) should be done, and I am by
no means sure I will go through this, not in this letter anyhow.

\margin{\unsure{ta.}} \hand{10.} Now in the long last the explicit description I promised of $\underline{C}_{n+1}$ \deleted{in} \inserted{in}
terms of $\underline{C}_n$ . As a matter of fact, I have a handful to propose ! One \margin{\unsure{n}}
choice, about the widest I would think of, is: for every pair (f,g)
in $\underline{C}_n$ satisfying condition (1) of (A), add one \xout{xxxxx} new arrow h. To
avoid set-theoretic difficulties though, we better \xout{xxxxxx} first modify
the definition of $\underline{C}_o$ so that the set of its objects should be in the
universe we are working in, preferably even it be countable. Or else
and more reasonably, we will pick one h for every isomorphism class of
situations (f,g) in $\underline{C}_n$ . Another restriction\hand{,} to avoid too much redundan-%
cy - this was the first definition actually that flipped to my mind the
day before yesterday - is to add a \uline{new} h only when there is no ``old''
one, namely in $\underline{C}_n$ , serving the same purpose. Then it came to my mind
that there is a lot of redundancy still\hand{;} \inserted[typed]{thus there would be already \deleted{in}}
\margin{30}
